\documentclass[10pt,a4paper]{amsart}
\usepackage{amsmath,amssymb,amsthm,graphicx}
% For dealing with references we use the comment environment
\usepackage{verbatim}
\newenvironment{reference}{\comment}{\endcomment}
%\newenvironment{reference}{}{}
\newenvironment{slogan}{\comment}{\endcomment}
\newenvironment{history}{\comment}{\endcomment}
% For commutative diagrams we use Xy-pic
\usepackage[all]{xy}

% We use 2cell for 2-commutative diagrams.
\xyoption{2cell}
\UseAllTwocells
\theoremstyle{plain}
\newtheorem{theorem}[subsection]{Theorem}
\newtheorem{proposition}[subsection]{Proposition}
\newtheorem{lemma}[subsection]{Lemma}

\theoremstyle{definition}
\newtheorem{definition}[subsection]{Definition}
\newtheorem{example}[subsection]{Example}
\newtheorem{exercise}[subsection]{Exercise}
\newtheorem{situation}[subsection]{Situation}

\theoremstyle{remark}
\newtheorem{remark}[subsection]{Remark}
\newtheorem{remarks}[subsection]{Remarks}

\numberwithin{equation}{subsection}
% Macros
%
\def\lim{\mathop{\mathrm{lim}}\nolimits}
\def\colim{\mathop{\mathrm{colim}}\nolimits}
\def\Spec{\mathop{\mathrm{Spec}}}
\def\Hom{\mathop{\mathrm{Hom}}\nolimits}
\def\Ext{\mathop{\mathrm{Ext}}\nolimits}
\def\SheafHom{\mathop{\mathcal{H}\!\mathit{om}}\nolimits}
\def\SheafExt{\mathop{\mathcal{E}\!\mathit{xt}}\nolimits}
\def\Sch{\mathit{Sch}}
\def\Mor{\mathop{\mathrm{Mor}}\nolimits}
\def\Ob{\mathop{\mathrm{Ob}}\nolimits}
\def\Sh{\mathop{\mathit{Sh}}\nolimits}
\def\NL{\mathop{N\!L}\nolimits}
\def\CH{\mathop{\mathrm{CH}}\nolimits}
\def\proetale{{pro\text{-}\acute{e}tale}}
\def\etale{{\acute{e}tale}}
\def\QCoh{\mathit{QCoh}}
\def\Ker{\mathop{\mathrm{Ker}}}
\def\Im{\mathop{\mathrm{Im}}}
\def\Coker{\mathop{\mathrm{Coker}}}
\def\Coim{\mathop{\mathrm{Coim}}}

% Boxtimes
%
\DeclareMathSymbol{\boxtimes}{\mathbin}{AMSa}{"02}

%
% Macros for moduli stacks/spaces
%
\def\QCohstack{\mathcal{QC}\!\mathit{oh}}
\def\Cohstack{\mathcal{C}\!\mathit{oh}}
\def\Spacesstack{\mathcal{S}\!\mathit{paces}}
\def\Quotfunctor{\mathrm{Quot}}
\def\Hilbfunctor{\mathrm{Hilb}}
\def\Curvesstack{\mathcal{C}\!\mathit{urves}}
\def\Polarizedstack{\mathcal{P}\!\mathit{olarized}}
\def\Complexesstack{\mathcal{C}\!\mathit{omplexes}}
% \Pic is the operator that assigns to X its picard group, usage \Pic(X)
% \Picardstack_{X/B} denotes the Picard stack of X over B
% \Picardfunctor_{X/B} denotes the Picard functor of X over B
\def\Pic{\mathop{\mathrm{Pic}}\nolimits}
\def\Picardstack{\mathcal{P}\!\mathit{ic}}
\def\Picardfunctor{\mathrm{Pic}}
\def\Deformationcategory{\mathcal{D}\!\mathit{ef}}
\usepackage{cleveref}
\crefname{originallemma}{original lemma}{original lemmas}
\Crefname{originallemma}{Original lemma}{Original lemmas}
\makeatletter
\newlabel{fields-lemma-separable-first}{{030K}{source}{}{sourcepointer-1}{}}
\newlabel{fields-lemma-separable-first@cref}{{[originallemma][0][]030K}{[1][0][]source}{}{}{}}
\makeatother
\crefname{originalsection}{original section}{original sections}
\Crefname{originalsection}{Original section}{Original sections}
\makeatletter
\newlabel{fields-section-separable-extensions}{{09GZ}{source}{}{sourcepointer-2}{}}
\newlabel{fields-section-separable-extensions@cref}{{[originalsection][0][]09GZ}{[1][0][]source}{}{}{}}
\makeatother
\crefname{originallemma}{original lemma}{original lemmas}
\Crefname{originallemma}{Original lemma}{Original lemmas}
\makeatletter
\newlabel{fields-lemma-count-embeddings}{{09H7}{source}{}{sourcepointer-3}{}}
\newlabel{fields-lemma-count-embeddings@cref}{{[originallemma][0][]09H7}{[1][0][]source}{}{}{}}
\makeatother
\title{Lemma powers field in algebra.tex}
\author{The Stacks Project Authors}
\date{stacks-project:a04446e57ec1fbc252a871afcec7752fb2807b14; GNU Free Documentation License 1.2 (as published in the Stacks Project COPYING)}
\begin{document}
\maketitle
\noindent\textit{Editorial note.} All mathematical statement and proof text is reproduced verbatim from the author's chapter algebra.tex of the Stacks Project, fixed at commit a04446e57ec1fbc252a871afcec7752fb2807b14. The pinned preamble is supplied as its own file; the wrapper class, the xr-hyper wiring and the \textbackslash{}externaldocument list are omitted because the assembler supplies the class.\par\medskip
\noindent\textit{Editorial note.} Numbering inside this standalone excerpt is editorial, not the full Stacks book print numbering.\par\medskip
\noindent\textit{Editorial note.} The candidate is from the extractor's stacks-candidates line, not from the earlier judged initial\_corpus line; the H2 judgement and its reason below were made in this module by reading the author's own statement and complete proof in the fixed checkout. Original candidate key: stacks:algebra.tex:lemma-powers-field.\par\medskip
\noindent\textit{Editorial note.} Reference labels in the proof resolve to navigation pointers to the original author source; these pointers do not reproduce proofs of the prerequisites.\par\medskip
\noindent\textit{Allowed standard prerequisites (editorial).}\begin{itemize}
\item Prerequisites named by the author inside the counted span. The counted span is reproduced verbatim, including the author's own citations. Its prerequisites are the earlier results of the same fixed Stacks Project version that the author cites there and in the chapter's lead-in.
\end{itemize}\medskip
\ifdefined\hypertarget\hypertarget{sourcepointer-1}{}\fi
\noindent\textit{Source reference.} Original lemma 030K: sources/fields.tex, source lines 1704--1704. Navigation pointer used in the counted span; the span contains the pointer only and the cited result is not reproduced inside it. Registered as a non-load-bearing navigation stub in the established form for Stacks material. No bibliography citation is involved.\par\medskip
\ifdefined\hypertarget\hypertarget{sourcepointer-2}{}\fi
\noindent\textit{Source reference.} Original section 09GZ: sources/fields.tex, source lines 1124--1124. Navigation pointer used in the counted span; the span contains the pointer only and the cited result is not reproduced inside it. Registered as a non-load-bearing navigation stub in the established form for Stacks material. No bibliography citation is involved.\par\medskip
\ifdefined\hypertarget\hypertarget{sourcepointer-3}{}\fi
\noindent\textit{Source reference.} Original lemma 09H7: sources/fields.tex, source lines 1291--1291. Navigation pointer used in the counted span; the span contains the pointer only and the cited result is not reproduced inside it. Registered as a non-load-bearing navigation stub in the established form for Stacks material. No bibliography citation is involved.\par\medskip
\begin{lemma}
\label{lemma-powers-field}
\begin{reference}
\cite[Lemma 3.1.6]{Alper-adequate}
\end{reference}
Let $k'/k$ be a field extension. The following are equivalent
\begin{enumerate}
\item for each $x \in k'$ there exists an $n > 0$ such that $x^n \in k$, and
\item $k' = k$ or $k$ and $k'$ have characteristic $p > 0$ and
either $k'/k$ is a purely inseparable extension or
$k$ and $k'$ are algebraic extensions of $\mathbf{F}_p$.
\end{enumerate}
\end{lemma}

\begin{proof}
Observe that each of the possibilities listed in (2) satisfies (1).
Thus we assume $k'/k$ satisfies (1) and we prove that we are in
one of the cases of (2). Discarding the case $k = k'$ we may assume
$k' \not = k$. It is clear that $k'/k$ is algebraic.
Hence we may assume that $k'/k$ is a nontrivial finite extension.
Let $k'/k'_{sep}/k$ be the separable subextension
found in Fields, Lemma \ref{fields-lemma-separable-first}.
We have to show that $k = k'_{sep}$ or that $k$ is an algebraic over
$\mathbf{F}_p$. Thus we may assume that $k'/k$
is a nontrivial finite separable extension and we have to show
$k$ is algebraic over $\mathbf{F}_p$.

\medskip\noindent
Pick $x \in k'$, $x \not \in k$. Pick $n, m > 0$ such that
$x^n \in k$ and $(x + 1)^m \in k$. Let $\overline{k}$ be an
algebraic closure of $k$. We can choose embeddings
$\sigma, \tau : k' \to \overline{k}$ with $\sigma(x) \not = \tau(x)$.
This follows from the discussion in
Fields, Section \ref{fields-section-separable-extensions}
(more precisely, after replacing $k'$ by the $k$-extension
generated by $x$ it follows from
Fields, Lemma \ref{fields-lemma-count-embeddings}).
Then we see that $\sigma(x) = \zeta \tau(x)$ for some
$n$th root of unity $\zeta$ in $\overline{k}$.
Similarly, we see that $\sigma(x + 1) = \zeta' \tau(x + 1)$
for some $m$th root of unity $\zeta' \in \overline{k}$.
Since $\sigma(x + 1) \not = \tau(x + 1)$ we see $\zeta' \not = 1$.
Then
$$
\zeta' (\tau(x) + 1) =
\zeta' \tau(x + 1) =
\sigma(x + 1) =
\sigma(x) + 1 =
\zeta \tau(x) + 1
$$
implies that
$$
\tau(x) (\zeta' - \zeta) = 1 - \zeta'
$$
hence $\zeta' \not = \zeta$ and
$$
\tau(x) = (1 - \zeta')/(\zeta' - \zeta)
$$
Hence every element of $k'$ which is not in $k$ is algebraic over the prime
subfield. Since $k'$ is generated over the prime subfield by the elements
of $k'$ which are not in $k$, we conclude that $k'$ (and hence $k$)
is algebraic over the prime subfield.

\medskip\noindent
Finally, if the characteristic of $k$ is $0$, the above leads to a
contradiction as follows (we encourage the reader to find their own proof).
For every rational number $y$ we similarly get a root of unity
$\zeta_y$ such that $\sigma(x + y) = \zeta_y\tau(x + y)$.
Then we find
$$
\zeta \tau(x) + y = \zeta_y(\tau(x) + y)
$$
and by our formula for $\tau(x)$ above we conclude
$\zeta_y \in \mathbf{Q}(\zeta, \zeta')$. Since the number field
$\mathbf{Q}(\zeta, \zeta')$ contains only a finite number of roots of
unity we find two distinct rational numbers $y, y'$ with
$\zeta_y = \zeta_{y'}$. Then we conclude that
$$
y - y' =
\sigma(x + y) - \sigma(x + y') =
\zeta_y(\tau(x + y)) - \zeta_{y'}\tau(x + y') = \zeta_y(y - y')
$$
which implies $\zeta_y = 1$ a contradiction.
\end{proof}

\end{document}
